% GENERATED FILE. Edit canonical lessons, then run npm run generate:books.
% canonical-source-hash: fnv1a64:b94d5a34dc246f64
% canonical-lessons: GE-C472-sehnen, GE-C472-sehnsucht, GE-C472-einsamkeit, GE-C472-heimweh, GE-C472-vermissen

\chapter{Sehnsucht und Heimweh}
\label{ch:ge-sehnsucht-und-heimweh}
\hlchaptermodality{472}

\begin{chapteropening}
\textbf{By the end of this chapter:} \emph{I can talk about longing, loneliness and missing people.}

Complete the last lesson of chapter 472: I can talk about longing, loneliness and missing people.
\end{chapteropening}

\section[sehnen]{sehnen — to long (sich sehnen nach)}

\label{lesson:GE-C472-sehnen}



\begin{quote}
Before the new one: say the German for to comfort, then the German for pity.
\end{quote}

\subsection*{You'll want to know: sehnen}
\textbf{sehnen} — \textquotedblleft{}to long (sich sehnen nach)\textquotedblright{}.

\begin{quote}
Im Winter sehne ich mich nach der Sonne.
\end{quote}

\begin{quote}
\emph{In winter I long for the sun.}
\end{quote}

Reflexive: \emph{sich nach etwas sehnen}.

\begin{cousinweb}
\emph{sehnen} goes back to Middle High German \emph{senen} \textquotedblleft{}to long, to pine\textquotedblright{}; its older origin is not settled, and EtymDB lists the German verb without a source.
\end{cousinweb}

\begin{grammarlens}[title={What you've built}]
One more word for telling, explaining and describing at B1.
\end{grammarlens}

\subsection*{Your turn}
\begin{itemize}
\raggedright
  \item \emph{Say it:} \emph{sehnen}
  \item \emph{Say it:} \emph{sehnen}, once more
  \item \emph{Say it:} a sentence of your own with \emph{sehnen}
  \item \emph{Recall:} say the German for to comfort, then the German for pity, then say \emph{sehnen} again
\end{itemize}

\begin{tcolorbox}[breakable,colback=teal!4,colframe=teal!35!black,title={Before you move on}]
What does sehnen mean? (\textquotedblleft{}To long (sich sehnen nach)\textquotedblright{}.) Say it once more.
\end{tcolorbox}
\section[die Sehnsucht]{die Sehnsucht — longing, yearning}

\label{lesson:GE-C472-sehnsucht}



\begin{quote}
Before the new one: say the German for pity, then the German for to long.
\end{quote}

\subsection*{You'll want to know: die Sehnsucht}
\textbf{die Sehnsucht} — \textquotedblleft{}longing, yearning\textquotedblright{}.

\begin{quote}
Er hatte große Sehnsucht nach seiner Familie.
\end{quote}

\begin{quote}
\emph{He longed for his family.}
\end{quote}

Plural: \emph{die Sehnsüchte}. \emph{Sehnsucht nach} + dative.

\begin{cousinweb}
\emph{Sehnsucht} is built on \emph{sehnen} \textquotedblleft{}to long\textquotedblright{} and \emph{Sucht} \textquotedblleft{}sickness, craving\textquotedblright{} (EtymDB), which is inherited from Old High German \emph{suht}, from Proto-Germanic \emph{suhtiz}.
\end{cousinweb}

\begin{grammarlens}[title={What you've built}]
One more word for telling, explaining and describing at B1.
\end{grammarlens}

\subsection*{Your turn}
\begin{itemize}
\raggedright
  \item \emph{Say it:} \emph{die Sehnsucht}
  \item \emph{Say it:} \emph{die Sehnsucht}, once more
  \item \emph{Say it:} a sentence of your own with \emph{die Sehnsucht}
  \item \emph{Recall:} say the German for pity, then the German for to long, then say \emph{die Sehnsucht} again
\end{itemize}

\begin{tcolorbox}[breakable,colback=teal!4,colframe=teal!35!black,title={Before you move on}]
What does die Sehnsucht mean? (\textquotedblleft{}Longing, yearning\textquotedblright{}.) Say it once more.
\end{tcolorbox}
\section[die Einsamkeit]{die Einsamkeit — loneliness; solitude}

\label{lesson:GE-C472-einsamkeit}



\begin{quote}
Before the new one: say the German for to long, then the German for longing.
\end{quote}

\subsection*{You'll want to know: die Einsamkeit}
\textbf{die Einsamkeit} — \textquotedblleft{}loneliness; solitude\textquotedblright{}.

\begin{quote}
Nach dem Umzug litt sie unter Einsamkeit.
\end{quote}

\begin{quote}
\emph{After the move she suffered from loneliness.}
\end{quote}

No plural. Built on \emph{einsam} — lonely.

\begin{cousinweb}
\emph{Einsamkeit} is built on \emph{einsam} \textquotedblleft{}lonely\textquotedblright{} with \emph{-keit} (EtymDB lists both); \emph{einsam} joins \emph{ein} \textquotedblleft{}one\textquotedblright{}, from Proto-Germanic \emph{ainaz}, to the ending \emph{-sam}.
\end{cousinweb}

\begin{grammarlens}[title={What you've built}]
One more word for telling, explaining and describing at B1.
\end{grammarlens}

\subsection*{Your turn}
\begin{itemize}
\raggedright
  \item \emph{Say it:} \emph{die Einsamkeit}
  \item \emph{Say it:} \emph{die Einsamkeit}, once more
  \item \emph{Say it:} a sentence of your own with \emph{die Einsamkeit}
  \item \emph{Recall:} say the German for to long, then the German for longing, then say \emph{die Einsamkeit} again
\end{itemize}

\begin{tcolorbox}[breakable,colback=teal!4,colframe=teal!35!black,title={Before you move on}]
What does die Einsamkeit mean? (\textquotedblleft{}Loneliness; solitude\textquotedblright{}.) Say it once more.
\end{tcolorbox}
\section[das Heimweh]{das Heimweh — homesickness}

\label{lesson:GE-C472-heimweh}



\begin{quote}
Before the new one: say the German for longing, then the German for loneliness.
\end{quote}

\subsection*{You'll want to know: das Heimweh}
\textbf{das Heimweh} — \textquotedblleft{}homesickness\textquotedblright{}.

\begin{quote}
In der ersten Woche im Ferienlager hatte er Heimweh.
\end{quote}

\begin{quote}
\emph{In the first week at summer camp he was homesick.}
\end{quote}

No plural. The opposite is \emph{das Fernweh} — the urge to travel.

\begin{cousinweb}
\emph{Heimweh} is built on \emph{Heim} \textquotedblleft{}home\textquotedblright{} and \emph{Weh} \textquotedblleft{}pain, woe\textquotedblright{}, which goes back to Middle High German \emph{we} (EtymDB); the word is usually said to have spread from Swiss German.
\end{cousinweb}

\begin{grammarlens}[title={What you've built}]
One more word for telling, explaining and describing at B1.
\end{grammarlens}

\subsection*{Your turn}
\begin{itemize}
\raggedright
  \item \emph{Say it:} \emph{das Heimweh}
  \item \emph{Say it:} \emph{das Heimweh}, once more
  \item \emph{Say it:} a sentence of your own with \emph{das Heimweh}
  \item \emph{Recall:} say the German for longing, then the German for loneliness, then say \emph{das Heimweh} again
\end{itemize}

\begin{tcolorbox}[breakable,colback=teal!4,colframe=teal!35!black,title={Before you move on}]
What does das Heimweh mean? (\textquotedblleft{}Homesickness\textquotedblright{}.) Say it once more.
\end{tcolorbox}
\section[vermissen]{vermissen — to miss (someone or something)}

\label{lesson:GE-C472-vermissen}



\begin{quote}
Before the new one: say the German for loneliness, then the German for homesickness.
\end{quote}

\subsection*{You'll want to know: vermissen}
\textbf{vermissen} — \textquotedblleft{}to miss (someone or something)\textquotedblright{}.

\begin{quote}
Ich vermisse meine Freunde aus der Schule.
\end{quote}

\begin{quote}
\emph{I miss my friends from school.}
\end{quote}

Also to find something missing: \emph{Seit gestern vermisse ich meinen Schlüssel.}

\begin{cousinweb}
\emph{vermissen} is built on \emph{ver-} and the older verb \emph{missen} \textquotedblleft{}to lack, to do without\textquotedblright{} (EtymDB), a Germanic word; English \emph{miss} is a cousin.
\end{cousinweb}

\begin{grammarlens}[title={What you've built}]
One more word for telling, explaining and describing at B1.
\end{grammarlens}

\subsection*{Your turn}
\begin{itemize}
\raggedright
  \item \emph{Say it:} \emph{vermissen}
  \item \emph{Say it:} \emph{vermissen}, once more
  \item \emph{Say it:} a sentence of your own with \emph{vermissen}
  \item \emph{Recall:} say the German for loneliness, then the German for homesickness, then say \emph{vermissen} again
\end{itemize}

\begin{tcolorbox}[breakable,colback=teal!4,colframe=teal!35!black,title={Before you move on}]
What does vermissen mean? (\textquotedblleft{}To miss (someone or something)\textquotedblright{}.) Say it once more.
\end{tcolorbox}
